\begin{table*}[t]
\centering
\footnotesize
\setlength{\tabcolsep}{5pt}
\begin{tabular}{@{}p{0.2\textwidth}p{0.75\textwidth}@{}}
\toprule
\textbf{Term} & \textbf{Gloss} \\
\midrule
\emph{Falllösung} & Written case solution --- the end-to-end work product that is graded. \\
\emph{Gutachtenstil} & Appraisal style: the conclusion follows the reasoning; the register graded here. \\
\emph{Urteilsstil} & Judgment style: result first, reasoning after; used in judicial decisions. \\
\emph{Obersatz} & Step 1 of the scaffold: the norm hypothesis to be examined. \\
\emph{Definition} & Step 2: interpretation of the norm's elements. \\
\emph{Subsumtion} & Step 3: anchoring the concrete facts under the norm's elements; broadly, the entire norm-application method. \\
\emph{Ergebnis} & Step 4: the resulting legal consequence. \\
\emph{Musterlösung} & Expert reference solution against which submissions are graded. \\
\emph{Staatsexamen} & State law examination; gate to all regulated legal professions in Germany. \\
\emph{ausreichend} & ``Sufficient'' --- lowest passing grade (4--6 of 18 points). \\
\emph{vollbefriedigend} & ``Fully satisfactory'' (10--12 points); customary distinction tier (\emph{Prädikat}). \\
\emph{Referendar} & Legal trainee in the two-year practical stage between the state examinations. \\
\emph{Volljurist} & Fully qualified jurist (passed both state examinations). \\
\emph{Grundprinzipien} & Core doctrinal principles; German name of the Doctrinal Principles corpus. \\
\emph{Bereich(e)} & Legal area(s): civil, public, criminal law. \\
Ja/Nein & Yes/no decision label on Doctrinal Principles items. \\
\emph{Ergebnisrichtigkeit} & Rubric dimension: result correctness. \\
\emph{Vollständigkeit} & Rubric dimension: issue identification / completeness. \\
\emph{Rechtsgrundlagen} & Rubric dimension: legal grounding (norm citation). \\
\emph{Rechtskenntnis} & Rubric dimension: doctrinal knowledge. \\
\emph{Schwerpunktsetzung} & Rubric dimension: problem depth at the case's focal issues. \\
\emph{Methodischer Stil} & Rubric dimension: methodological structure (scaffold execution). \\
\emph{Gliederung} & Rubric dimension: organization. \\
\emph{Sprache} & Rubric dimension: terminology and language. \\
\emph{Formalia} & Rubric dimension: formal correctness. \\
\bottomrule
\end{tabular}
\caption{German terms used in the paper, with English glosses.}
\label{tbl-glossary}
\end{table*}
